% Part: second-order-logic
% Chapter: syntax-and-semantics
% Section: satisfaction

\documentclass[../../../include/open-logic-section]{subfiles}

\begin{document}

\olfileid[tr]{sol}{syn}{sat}

\olsection{Sağlama}

\begin{explain}
İkinci dereceden !!{formula}s için sağlama bağıntısı $\Sat{M}{!A}[s]$'yi
tanımlamak üzere, tanımları ikinci dereceden !!{variable}s öğelerini
kapsayacak biçimde genişletmemiz gerekir. !!a{structure} kavramı ikinci
dereceden mantıkta da birinci dereceden mantıktakiyle aynıdır. Yalnızca
değişken atamaları~$s$ bakımından bir fark vardır: bunlar artık yalnızca
birinci dereceden !!{variable}s öğelerine değil, ikinci dereceden
!!{variable}s öğelerine de değer vermelidir.
\end{explain}

\begin{defn}[Değişken ataması]
!!a{structure}~$\Struct{M}$ için bir \emph{değişken ataması}~$s$, her
\begin{enumerate}
\item nesne !!{variable}~$\Obj{v_i}$ öğesini~$\Domain M$'nin bir öğesine,
  yani $s(\Obj{v_i}) \in \Domain{M}$ olacak biçimde;
\item $n$-li bağıntı değişkeni~$\Obj{V_i^n}$ öğesini~$\Domain{M}$ üzerinde
  $n$-li bir bağıntıya, yani $s(\Obj{V_i^n}) \subseteq \Domain{M}^n$
  olacak biçimde;
\item $n$-li fonksiyon değişkeni~$\Obj{u_i^n}$ öğesini $\Domain{M}$'den
  $\Domain{M}$'ye $n$-li bir fonksiyona, yani
  $s(\Obj{u_i^n})\colon \Domain{M}^n \to \Domain{M}$ olacak biçimde
\end{enumerate}
gönderen bir fonksiyondur.
\end{defn}

\begin{explain}
!!^a{structure}, !!a{value} atayarak her !!{constant} ve !!{function}
öğesini, bir bağıntı atayarak da her yüklemi yorumlar; ikinci dereceden bir değişken ataması ise her
nesne, bağıntı ve fonksiyon değişkenine sırasıyla bir nesne, bağıntı ve
fonksiyon atar. İkisi birlikte her terime bir değer atamamızı sağlar.
\end{explain}

\begin{defn}[Bir terimin \usetoken{S}{value}]
$t$,~$\Lang L$ dilinin bir terimi, $\Struct M$,~$\Lang L$ için
!!a{structure} ve $s$ !!a{variable} ataması bu~$\Struct M$ yapısına ait ise
\emph{!!{value}}~$\Value{t}{M}[s]$, birinci dereceden terimlerdeki gibi,
aşağıdaki ek maddeyle tanımlanır:
\begin{quote}
\indcase{t}{\Atom{u}{t_1, \ldots, t_n}}{
\[
\Value{\indfrm}{M}[s] = s(u)(\Value{t_1}{M}[s], \ldots,
\Value{t_n}{M}[s]).
\]}
\end{quote}
\end{defn}

\begin{defn}[$x$-değişkesi]
$s$ bir !!a{variable} ataması ve bu atama !!a{structure}~$\Struct M$'ye ait olsun.
Yine~$\Struct M$ için olan ve $s$'den en çok $x$'e atadığı değer bakımından
ayrılan her $s'$ !!{variable}
atamasına~$s$'nin bir \emph{$x$-değişkesi} denir. $s'$, $s$'nin
bir $x$-değişkesiyse $\varAssign{s'}{s}{x}$ yazarız. (İkinci dereceden
değişkenler $X$ ya da $u$ için de benzer biçimde.)
\end{defn}

\begin{defn}
  $s$ bir !!a{variable} ataması, bu atama !!a{structure}~$\Struct M$'ye ait ve
  $m \in \Domain{M}$ ise~$\Subst{s}{m}{x}$ ataması,
  \[\Subst{s}{m}{y} = \begin{cases}
    m & \text{eğer } y \ident x\\
    s(y) & \text{aksi hâlde},
  \end{cases}\]
  ile tanımlanan !!{variable} atamasıdır.
  $X$, $n$-li bir bağıntı !!{variable} ve $M \subseteq \Domain{M}^n$ ise
  $\Subst{s}{M}{X}$,
  \[\Subst{s}{M}{y} = \begin{cases}
    M & \text{eğer } y \ident X\\
    s(y) & \text{aksi hâlde}.
  \end{cases}\]
  ile tanımlanan !!{variable} atamasıdır.
  $u$, $n$-li bir fonksiyon !!{variable} ve
  $f\colon \Domain{M}^n \to \Domain{M}$ ise $\Subst{s}{f}{u}$,
  \[\Subst{s}{f}{y} = \begin{cases}
    f & \text{eğer } y \ident u\\
    s(y) & \text{aksi hâlde}.
  \end{cases}\]
  ile tanımlanan !!{variable} atamasıdır. Her durumda $y$, birinci
  ya da ikinci dereceden herhangi bir !!{variable} olabilir.
\end{defn}

\begin{defn}[Sağlama]
İkinci dereceden !!{formula}s~$!A$ için sağlama tanımı,
\olref[fol][syn][sat]{defn:satisfaction} tanımına şu maddelerin eklenmesiyle
elde edilir:
\begin{enumerate}
\item \indcase{!A}{\Atom{X^n}{t_1, \dots, t_n}}{$\Sat{M}{\indfrm}[s]$
  ancak ve ancak $\langle \Value{t_1}{M}[s], \dots,
  \Value{t_n}{M}[s] \rangle \in s(X^n)$ ise.}
\tagitem{prvAll}{%
  \indcase{!A}{\lforall[X][!B]}{$\Sat{M}{\indfrm}[s]$ ancak ve ancak her
  $M \subseteq \Domain{M}^n$ için
  $\Sat{M}{!B}[\Subst{s}{M}{X}]$ ise.}}{}

\tagitem{prvEx}{%
  \indcase{!A}{\lexists[X][!B]}{$\Sat{M}{\indfrm}[s]$ ancak ve ancak
  $\Sat{M}{!B}[\Subst{s}{M}{X}]$ olacak en az bir
  $M \subseteq \Domain{M}^n$ varsa.}}{}

\tagitem{prvAll}{%
  \indcase{!A}{\lforall[u][!B]}{$\Sat{M}{\indfrm}[s]$ ancak ve ancak her
  $f\colon \Domain{M}^n \to \Domain{M}$ için
  $\Sat{M}{!B}[\Subst{s}{f}{u}]$ ise.}}{}

\tagitem{prvEx}{%
  \indcase{!A}{\lexists[u][!B]}{$\Sat{M}{\indfrm}[s]$ ancak ve ancak
  $\Sat{M}{!B}[\Subst{s}{f}{u}]$ olacak en az bir
  $f\colon \Domain{M}^n \to \Domain{M}$ varsa.}}{}
\end{enumerate}
\end{defn}

\begin{ex}
  $\lforall[z][(\Atom{X}{z} \liff \lnot \Atom{Y}{z})]$ !!{formula}
  ifadesini ele alalım. İkinci dereceden niceleyici içermez; fakat burada
  birli oldukları kabul edilen ikinci dereceden $X$ ve~$Y$ !!{variable}s
  öğelerini içerir. Buna karşılık gelen birinci dereceden !!{sentence}
  $\lforall[z][(\Atom{P}{z} \liff \lnot \Atom{R}{z})]$, $P$'nin yorumu
  kapsamına giren her şeyin~$R$'nin yorumu kapsamına girmediğini, bunun
  tersinin de geçerli olduğunu söyler. !!a{structure} içinde
  !!a{predicate}~$P$'nin yorumu~$\Assign{P}{M}$ ile verilir. Ancak $X$
  ve~$Y$ gibi ikinci dereceden !!{variable}s öğelerinin yorumunu
  !!{structure} değil, !!a{variable} ataması verir. İkinci dereceden
  !!{formula} bir !!a{sentence} olmadığından (serbest $X$ ve~$Y$
  !!{variable}s öğelerini içerir), ancak !!a{structure}~$\Struct{M}$ ile
  birlikte !!a{variable} ataması~$s$'ye göre sağlanır.

  $s(X)$'in !!{element}s öğeleri $s(Y)$'nin !!{element}s öğeleri değilse
  ve bunun tersi de geçerliyse, yani ancak ve ancak
  $s(Y) = \Domain{M} \setminus s(X)$ ise
  $\Sat{M}{\lforall[z][(\Atom{X}{z} \liff \lnot \Atom{Y}{z})]}[s]$ olur.
  Örneğin $\Domain{M} = \{1, 2, 3\}$ olsun. Hiçbir !!{predicate},
  !!{function} ya da !!{constant} söz konusu olmadığından,
  $\Struct{M}$'nin yalnızca evreni önemlidir. Şimdi
  $s_1(X) = \{1, 2\}$ ve $s_1(Y) = \{3\}$ için
  $\Sat{M}{\lforall[z][(\Atom{X}{z}
      \liff \lnot \Atom{Y}{z})]}[s_1]$ olur.

  Buna karşılık $s_2(X) = \{1, 2\}$ ve $s_2(Y) = \{2, 3\}$ ise
  $\Sat/{M}{\lforall[z][(\Atom{X}{z} \liff \lnot
      \Atom{Y}{z})]}[s_2]$ olur. Çünkü
  $\Sat{M}{\Atom{X}{z}}[\Subst{s_2}{2}{z}]$ (zira
  $2 \in \Subst{s_2}{2}{z}(X)$), fakat
  $\Sat/{M}{\lnot \Atom{Y}{z}}[\Subst{s_2}{2}{z}]$ olur (zira aynı zamanda
  $2 \in \Subst{s_2}{2}{z}(Y)$).
\end{ex}

\begin{ex}
  $\Sat{M}{\lexists[Y][(\lexists[y][\Atom{Y}{y}] \land
  \lforall[z][(\Atom{X}{z} \liff \lnot \Atom{Y}{z})])]}[s]$ ancak ve
  ancak öyle bir $N \subseteq \Domain{M}$ vardır ki
  $\Sat{M}{(\lexists[y][\Atom{Y}{y}] \land \lforall[z][(\Atom{X}{z}
  \liff \lnot \Atom{Y}{z})])}[\Subst{s}{N}{Y}]$ olur. Bu da hem
  $N \neq \emptyset$ (dolayısıyla
  $\Sat{M}{\lexists[y][\Atom{Y}{y}]}[\Subst{s}{N}{Y}]$) hem de önceki
  örnekteki gibi $N = \Domain{M} \setminus s(X)$ olduğu her durumda
  geçerlidir. Başka bir deyişle,
  $\Sat{M}{\lexists[Y][(\lexists[y][\Atom{Y}{y}] \land
  \lforall[z][(\Atom{X}{z} \liff \lnot \Atom{Y}{z})])]}[s]$ ancak ve
  ancak $\Domain{M} \setminus s(X)$ boş değilse, yani
  $s(X) \neq \Domain{M}$ ise geçerlidir. Dolayısıyla !!{formula}, örneğin
  $\Domain{M} = \{1, 2, 3\}$ ve $s(X) = \{1, 2\}$ olduğunda sağlanır;
  fakat $s(X) = \{1, 2, 3\} = \Domain{M}$ olduğunda sağlanmaz.

  !!{formula}, $s(X) = \Domain{M}$ olduğunda sağlanmadığından,
  \[
  \lforall[X][\lexists[Y][(\lexists[y][\Atom{Y}{y}] \land
      \lforall[z][(\Atom{X}{z} \liff \lnot \Atom{Y}{z})])]]
  \]
  !!{sentence} ifadesi hiçbir zaman sağlanmaz: Her
  !!{structure}~$\Struct{M}$ için $s(X)=\Domain{M}$ ataması !!{sentence}
  ifadesini yanlış kılar. Öte yandan,
  \[
  \lexists[X][\lexists[Y][(\lexists[y][\Atom{Y}{y}] \land
      \lforall[z][(\Atom{X}{z} \liff \lnot \Atom{Y}{z})])]]
  \]
  tümcesi her~$s$ atamasına göre sağlanır; çünkü her zaman
  $M \subseteq \Domain{M}$ ve $M \neq \Domain{M}$ olacak bir alt küme
  (örneğin $M = \emptyset$) seçebiliriz.
\end{ex}

\begin{ex}
  İkinci dereceden
  !!{sentence}~$\lforall[X][\lforall[y][X(y)]]$, her $1$-li bağıntının,
  yani her özelliğin her nesne için geçerli olduğunu söyler. Bu açıkça
  hiçbir zaman doğru değildir; çünkü her~$\Struct{M}$ içinde
  $s(X)=\emptyset$ ve $s(y)=a\in\Domain{M}$ olan bir değişken ataması~$s$
  için $\Sat/{M}{X(y)}[s]$ olur. Bu, ikinci dereceden mantıkta
  $!A \lif \lforall[X][\lforall[y][X(y)]]$ ifadesinin $\lnot !A$ ile
  eşdeğer olduğu anlamına gelir; yani
  $\Sat{M}{!A \lif \lforall[X][\lforall[y][X(y)]]}$ ancak ve ancak
  $\Sat{M}{\lnot !A}$ ise geçerlidir. Başka bir deyişle, ikinci dereceden
  mantıkta $\lnot$ bağlacını $\lforall$ ve~$\lif$ aracılığıyla
  tanımlayabiliriz.
\end{ex}

\begin{prob}
  İkinci dereceden mantıkta $\lforall$ ve~$\lif$ ile öteki bağlaçların
  tanımlanabildiğini gösterin:
  \begin{enumerate}
    \item İkinci dereceden mantıkta $!A \land !B$ ifadesinin
    $\lforall[X][(!A \lif (!B \lif \lforall[x][X(x)])\lif
    \lforall[x][X(x)])]$ ile eşdeğer olduğunu kanıtlayın.
    \item Yalnızca $\lforall$ ve $\lif$ kullanan ve $!A \lor !B$ ile
    eşdeğer olan ikinci dereceden bir formül bulun.
  \end{enumerate}
\end{prob}

\end{document}
